\documentclass[12pt,a4paper]{amsart}

\usepackage{amsmath,amssymb,amsthm}
\usepackage{mathtools}
\usepackage[utf8]{inputenc}
\usepackage[ngerman]{babel}
\usepackage{hyperref}
\usepackage{geometry}
\geometry{margin=2.5cm}

% -------------------------------------------------------
% Theorem-Umgebungen
% -------------------------------------------------------
\newtheorem{theorem}{Satz}[section]
\newtheorem{lemma}[theorem]{Lemma}
\newtheorem{corollary}[theorem]{Korollar}
\newtheorem{proposition}[theorem]{Proposition}
\theoremstyle{definition}
\newtheorem{definition}[theorem]{Definition}
\newtheorem{remark}[theorem]{Bemerkung}

% -------------------------------------------------------
% Eigene Befehle
% -------------------------------------------------------
\newcommand{\N}{\mathbb{N}}
\newcommand{\Z}{\mathbb{Z}}
\newcommand{\R}{\mathbb{R}}
\newcommand{\C}{\mathbb{C}}
\DeclareMathOperator{\Re}{Re}
\DeclareMathOperator{\Im}{Im}

% -------------------------------------------------------
\begin{document}
% -------------------------------------------------------

\title{Die nicht-trivialen Nullstellen der Riemann'schen Zeta-Funktion:\\
       Z\"ahlung, Verteilung und das GUE-Gesetz}

\author{Michael Fuhrmann}
\address{Specialist Mathematics Project, Build 84}
\email{specialist-maths@localhost}
\date{\today}

\begin{abstract}
Wir untersuchen die nicht-trivialen Nullstellen der Riemann'schen Zeta-Funktion
$\zeta(s)$, d.\,h.\ die Nullstellen im kritischen Streifen $0 < \Re(s) < 1$.
Zentrales Ergebnis ist die \emph{Nullstellenz\"ahlformel}
\[
  N(T) = \frac{T}{2\pi}\log\frac{T}{2\pi} - \frac{T}{2\pi} + O(\log T),
\]
die die Anzahl der Nullstellen $\rho = \beta + i\gamma$ mit
$0 < \gamma \le T$ bestimmt.
Wir identifizieren die erste bekannte nicht-triviale Nullstelle
$\rho_1 = \tfrac{1}{2} + 14{,}134725\ldots\,i$ und diskutieren
Gram-Punkte sowie Grams Gesetz (einschlie\ss lich seiner Ausnahmen).
Das tiefe Montgomery--Odlyzko-Gesetz besagt, dass die lokale Statistik der
Nullstellenabst\"ande der \emph{GUE}-Verteilung (Gaussian Unitary Ensemble)
der Zufallsmatrizentheorie folgt.
Mehr als $10^{13}$ Nullstellen wurden numerisch auf der kritischen
Geraden $\Re(\rho) = \tfrac{1}{2}$ verifiziert.
\end{abstract}

\maketitle
\tableofcontents

% ================================================================
\section{Einleitung}
% ================================================================

Die nicht-trivialen Nullstellen von $\zeta(s)$ sind die Nullstellen
im kritischen Streifen $0 < \Re(s) < 1$.
Aus der Funktionalgleichung folgt: Ist $\rho$ eine Nullstelle, so auch
$1-\rho$. Aus dem Schwarz'schen Spiegelungsprinzip folgt: Ist $\rho$
eine Nullstelle, so auch $\bar\rho$.
Damit treten Nullstellen in Vierergruppen
$\{\rho, \bar\rho, 1-\rho, 1-\bar\rho\}$ auf, sofern nicht
$\Re(\rho) = \tfrac{1}{2}$ gilt -- dann bilden sie Zweierpaare
$\{\rho, \bar\rho\}$.

\begin{remark}
Alle bisher gefundenen Nullstellen erf\"ullen $\Re(\rho) = \tfrac{1}{2}$.
Die Riemann-Hypothese (RH) vermutet, dass dies f\"ur \emph{alle}
nicht-trivialen Nullstellen gilt; vgl.\ \cite{Riemann1859}.
\end{remark}

% ================================================================
\section{Z\"ahlung von Nullstellen: Die Formel \texorpdfstring{$N(T)$}{N(T)}}
% ================================================================

\begin{definition}[Nullstellenz\"ahlfunktion]
\label{def:NT}
F\"ur $T > 0$ sei
\[
  N(T) \;=\; \#\bigl\{\rho = \beta + i\gamma : \zeta(\rho) = 0,\;
  0 < \beta < 1,\; 0 < \gamma \le T\bigr\}
\]
die Anzahl der nicht-trivialen Nullstellen von $\zeta$ mit positivem
Imagin\"arteil $\le T$ (mit Vielfachheit gez\"ahlt).
\end{definition}

\begin{theorem}[Riemann--von-Mangoldt-Formel]
\label{thm:NT}
F\"ur $T > 0$, das kein Imagin\"arteil einer Nullstelle von $\zeta$ ist,
gilt:
\begin{equation}
  \label{eq:NT}
  N(T) \;=\; \frac{T}{2\pi}\log\frac{T}{2\pi} - \frac{T}{2\pi}
            + \frac{7}{8} + O\!\bigl(\log T\bigr).
\end{equation}
\end{theorem}

\begin{proof}[Beweisskizze]
Der Beweis verwendet das Argumentprinzip angewandt auf das Rechteck
mit Ecken $-1$, $2$, $2+iT$, $-1+iT$.
Nach dem Argumentprinzip gilt $N(T) = \tfrac{1}{2\pi}\Delta\arg\zeta(s)$
\"uber dieses Rechteck.
Den Hauptbeitrag liefern die Segmente bei $\Re(s) = 2$
(dort ist $\zeta$ durch ihre Dirichlet-Reihe gen\"ahert)
und bei $\Re(s) = 0$ (behandelt via Funktionalgleichung).
Die Stirling-Formel f\"ur $\Gamma$ liefert die f\"uhrenden Terme.
Die Konstante $7/8$ und das Fehlerglied $O(\log T)$ folgen aus
sorgf\"altiger Buchf\"uhrung; vgl.\ \cite{Davenport2000}, Kapitel~16.
\end{proof}

\begin{corollary}[Dichte der Nullstellen]
Der mittlere Abstand zwischen aufeinanderfolgenden Nullstellen bei
H\"ohe $T$ betr\"agt n\"aherungsweise $2\pi/\log T$.
Die Nullstellen werden also f\"ur $T \to \infty$ immer dichter.
\end{corollary}

\begin{proof}
Ableitung von~\eqref{eq:NT}: $N'(T) \approx \frac{1}{2\pi}\log\frac{T}{2\pi}$,
also mittlerer Abstand $\approx 2\pi/\log T$.
\end{proof}

% ================================================================
\section{Die ersten nicht-trivialen Nullstellen}
% ================================================================

\begin{theorem}[Erste nicht-triviale Nullstelle]
\label{thm:first_zero}
Der kleinste positive Imagin\"arteil einer nicht-trivialen Nullstelle von
$\zeta$ ist
\[
  \gamma_1 = 14{,}134725141734693\ldots
\]
Somit: $\rho_1 = \tfrac{1}{2} + 14{,}134725\ldots\,i$.
\end{theorem}

\begin{proof}[Verifikation]
Durch direkte numerische Berechnung mit der Riemann--Siegel-Formel hat
$Z(t)$ f\"ur $0 < t < 14{,}134\ldots$ keine Nullstelle und bei
$t = 14{,}134725\ldots$ eine solche.
$N(14) = 0$ und $N(15) = 1$ werden explizit verifiziert.
Die numerischen Werte sind in der LMFDB-Datenbank \cite{LMFDB}
mit hoher Pr\"azision gespeichert.
\end{proof}

\begin{remark}[Die ersten zehn Nullstellen]
Die Imagin\"arteile der ersten zehn nicht-trivialen Nullstellen auf der
kritischen Geraden sind:
\begin{align*}
  \gamma_1 &= 14{,}134725\ldots, &
  \gamma_2 &= 21{,}022040\ldots, &
  \gamma_3 &= 25{,}010858\ldots, \\
  \gamma_4 &= 30{,}424876\ldots, &
  \gamma_5 &= 32{,}935062\ldots, &
  \gamma_6 &= 37{,}586178\ldots, \\
  \gamma_7 &= 40{,}918720\ldots, &
  \gamma_8 &= 43{,}327073\ldots, &
  \gamma_9 &= 48{,}005150\ldots, \\
  \gamma_{10} &= 49{,}773832\ldots
\end{align*}
Alle liegen auf der kritischen Geraden, konsistent mit RH.
\end{remark}

% ================================================================
\section{Gram-Punkte und Grams Gesetz}
% ================================================================

\begin{definition}[Gram-Punkte]
\label{def:gram_points}
Ein \emph{Gram-Punkt} $g_n$ (f\"ur $n \ge -1$) ist die eindeutige L\"osung von
\[
  \theta(g_n) = n\pi,
\]
wobei $\theta(t) = \arg\Gamma\!\bigl(\tfrac{1}{4}+\tfrac{it}{2}\bigr)
- \tfrac{t}{2}\ln\pi$ die Riemann--Siegel-Theta-Funktion ist.
Der Gram-Abstand $[g_n, g_{n+1}]$ hat L\"ange $\approx 2\pi/\log(g_n/2\pi)$.
\end{definition}

\begin{theorem}[Grams Gesetz und seine Verletzungen]
\label{thm:gram}
\emph{Grams Gesetz} (eine Beobachtung, kein Satz) besagt, dass $Z(g_n)$
das Vorzeichen $(-1)^n$ tr\"agt, sodass jedes Gram-Intervall
$[g_n, g_{n+1}]$ genau eine Nullstelle von $Z(t)$ enth\"alt.

Dies scheitert erstmals bei $n = 126$: Das Intervall
$[g_{126}, g_{127}]$ enth\"alt keine Nullstelle, w\"ahrend
$[g_{127}, g_{128}]$ zwei enth\"alt.
\end{theorem}

\begin{proof}[Diskussion]
Grams Gesetz ist eine Heuristik, die in etwa $73\,\%$ der F\"alle gilt
(Hutchinson, 1925).
Die Verletzung h\"angt mit dem Verhalten von $Z(t)$ an Gram-Punkten zusammen,
wo Real- und Imagin\"arteil von $\zeta(\tfrac{1}{2}+it)$ sich gerade aufheben.
Eine vollst\"andige Analyse findet sich in \cite{Titchmarsh1986}, Kapitel~9.
\end{proof}

% ================================================================
\section{Das Montgomery--Odlyzko-Gesetz}
% ================================================================

Die bedeutendste Entdeckung \"uber die Nullstellenverteilung ist ihre
Verbindung zur Zufallsmatrizentheorie.

\begin{definition}[Normalisierte Null\-stellen\-abst\"ande]
\label{def:spacing}
Seien $\rho_n = \tfrac{1}{2} + i\gamma_n$ die nicht-trivialen Nullstellen mit
$0 < \gamma_1 \le \gamma_2 \le \ldots$\,.
Der \emph{normalisierte Abstand} zwischen aufeinanderfolgenden Nullstellen ist
\[
  \delta_n \;=\; (\gamma_{n+1} - \gamma_n) \cdot \frac{\log\gamma_n}{2\pi}.
\]
Diese Normierung stellt sicher, dass $\delta_n$ im Mittel gegen~$1$ strebt.
\end{definition}

\begin{theorem}[Montgomerys Paarkorrelationsvermutung, 1973]
\label{thm:montgomery}
Unter Annahme der Riemann-Hypothese vermutete Montgomery \cite{Montgomery1973}
f\"ur die Zweipunkt-Korrelationsfunktion:
\[
  \lim_{T \to \infty} \frac{1}{N(T)}
  \sum_{\substack{\gamma, \gamma' \le T \\ 0 < \gamma'-\gamma \le 2\pi u / \log T}}
  1
  \;=\;
  \int_0^u \left[1 - \left(\frac{\sin(\pi x)}{\pi x}\right)^2\right] dx.
\]
Dies ist genau die Paarkorrelationsfunktion der Eigenwerte aus dem
\emph{Gaussian Unitary Ensemble} (GUE) zuf\"alliger hermitescher Matrizen.
\end{theorem}

\begin{remark}[Die GUE-Verbindung]
Die GUE-Paarkorrelationsfunktion $1 - (\sin\pi x/\pi x)^2$ impliziert
\emph{Niveauabsto\ss ung}: Nullstellen meiden sehr kleine gegenseitige Abst\"ande.
Dieselbe Funktion beschreibt die Abstandsstatistik von Kernenergieniveaus.
Odlyzko \cite{Odlyzko1987} entdeckte dies 1987 numerisch.
Die \"Ubereinstimmung gilt heute als \"uberzeugendes numerisches Indiz f\"ur RH.
\end{remark}

% ================================================================
\section{Numerische Verifikation der Riemann-Hypothese}
% ================================================================

\begin{theorem}[Verifikation der RH bis zu gro\ss er H\"ohe]
\label{thm:rh_verified}
Die Riemann-Hypothese wurde numerisch f\"ur alle Nullstellen mit
Imagin\"arteil $0 < \gamma \le T_{\max}$ verifiziert:
\begin{itemize}
  \item Backlund (1914): $T = 200$ (79 Nullstellen)
  \item Turing (1953): $T \approx 1\,500$ (erste maschinelle Berechnung)
  \item Lune--te~Riele--Winter (1986): erste $1{,}5 \times 10^9$ Nullstellen
  \item Platt--Trudgian (2021) \cite{Platt2021}: alle Nullstellen unterhalb
    $T = 3 \times 10^{12}$
  \item Stand 2023: mehr als $10^{13}$ Nullstellen auf der kritischen Geraden
    verifiziert.
\end{itemize}
\end{theorem}

\begin{proof}[Methode]
Alle Verifikationen kombinieren:
\begin{enumerate}
  \item Die Riemann--Siegel-Formel zur effizienten Auswertung von $Z(t)$.
  \item Backlunds Methode: Vorzeichenwechsel von $Z(t)$ zwischen aufeinanderfolgenden
    Gram-Punkten z\"ahlen und mit $N(T)$ aus der Riemann--von-Mangoldt-Formel
    vergleichen.
  \item Sorgf\"altige Behandlung von \glqq Gram-Ausnahmen\grqq{}
    mit Turings Methode.
\end{enumerate}
Keine Nullstelle abseits der kritischen Geraden wurde je gefunden.
\end{proof}

% ================================================================
\section{Die Verallgemeinerte Riemann-Hypothese}
% ================================================================

\begin{definition}[GRH]
\label{def:grh}
Die \emph{Verallgemeinerte Riemann-Hypothese} (GRH) erstreckt RH auf alle
\emph{Dirichlet-L-Funktionen}:
\[
  L(s, \chi) = \sum_{n=1}^{\infty} \frac{\chi(n)}{n^s},
\]
wobei $\chi$ ein Dirichlet-Charakter modulo $q$ ist.
GRH behauptet, dass alle nicht-trivialen Nullstellen von $L(s,\chi)$
auf $\Re(s) = \tfrac{1}{2}$ liegen.
\end{definition}

\begin{remark}
GRH hat weitreichende Konsequenzen f\"ur Primzahlen in arithmetischen
Progressionen. Beispielsweise folgt: F\"ur jeden Modul $q$ enth\"alt jede
arithmetische Progression $a, a+q, a+2q,\ldots$ (mit $\gcd(a,q)=1$) eine
Primzahl $p \le C\,q^2\,(\log q)^2$.
\end{remark}

% ================================================================
\section{Schlussfolgerung}
% ================================================================

Die nicht-trivialen Nullstellen von $\zeta(s)$ zeigen eine reiche Struktur:
Die Z\"ahlformel~\eqref{eq:NT} bestimmt ihre makroskopische Verteilung,
die erste Nullstelle $\rho_1 \approx \tfrac{1}{2} + 14{,}13i$ wurde hochpr\"azise
lokalisiert, und die GUE-Statistik von Montgomery und Odlyzko offenbart
eine tiefe Verbindung zur Quantenmechanik.
Trotz mehr als $10^{13}$ verifizierten Nullstellen auf der kritischen
Geraden bleibt die Riemann-Hypothese eines der gro\ss ten ungel\"osten Probleme
der Mathematik.

% ================================================================
\begin{thebibliography}{10}

\bibitem{Riemann1859}
B.~Riemann,
\emph{\"Uber die Anzahl der Primzahlen unter einer gegebenen Gr\"o\ss e},
Monatsberichte der Berliner Akademie (1859), 671--680.

\bibitem{Davenport2000}
H.~Davenport,
\emph{Multiplicative Number Theory}, 3.~Aufl.\
(revidiert von H.~L.~Montgomery),
Springer, New York, 2000.

\bibitem{Titchmarsh1986}
E.~C.~Titchmarsh,
\emph{The Theory of the Riemann Zeta-Function}, 2.~Aufl.\
(revidiert von D.~R.~Heath-Brown),
Oxford University Press, 1986.

\bibitem{Montgomery1973}
H.~L.~Montgomery,
The pair correlation of zeros of the zeta function,
in: \emph{Analytic Number Theory}, Proc.\ Sympos.\ Pure Math.~XXIV,
Amer.\ Math.\ Soc., Providence, RI, 1973, S.~181--193.

\bibitem{Odlyzko1987}
A.~M.~Odlyzko,
On the distribution of spacings between zeros of the zeta function,
\emph{Math.\ Comp.}\ \textbf{48} (1987), Nr.~177, 273--308.

\bibitem{Platt2021}
D.~J.~Platt und T.~S.~Trudgian,
The Riemann hypothesis is true up to $3 \times 10^{12}$,
\emph{Bull.\ London Math.\ Soc.}\ \textbf{53} (2021), 792--797.

\bibitem{LMFDB}
The LMFDB Collaboration,
\emph{The $L$-functions and Modular Forms Database},
\url{https://www.lmfdb.org}, 2024.

\end{thebibliography}

\end{document}
